\documentclass[12pt,a4paper]{article}
\usepackage[margin=25mm, headheight=15pt, headsep=10pt]{geometry}
\usepackage{fontspec}
\usepackage[table]{xcolor} % Note: table option is required for rowcolors
\usepackage{titlesec}
\usepackage{tocloft}
\usepackage{fancyhdr}
\usepackage{booktabs}
\usepackage{tcolorbox}
\tcbuselibrary{skins, breakable}
\usepackage[colorlinks=true, linkcolor=emerald, urlcolor=emerald, bookmarks=true]{hyperref}

\definecolor{emerald}{RGB}{0,102,51}
\definecolor{darkred}{RGB}{139,0,0}
\definecolor{lightgray}{RGB}{245,245,245}

\usepackage[nil]{babel}
\babelprovide[import, main]{bengali}
\babelprovide[import]{arabic}
\babelfont{rm}[Renderer=HarfBuzz,Scale=1.0]{Noto Serif Bengali}
\babelfont{sf}[Renderer=HarfBuzz]{Noto Sans Bengali}
\babelfont{tt}[Renderer=HarfBuzz]{Noto Sans Bengali}
\babelfont[arabic]{rm}[Renderer=HarfBuzz,Scale=1.1]{Noto Naskh Arabic}

% Section styling
\titleformat{\section}{\Large\bfseries\color{emerald}}{\thesection.}{0.5em}{}
\titleformat{\subsection}{\large\bfseries\color{emerald!85!black}}{\thesubsection.}{0.5em}{}
\titleformat{\subsubsection}{\normalsize\bfseries\color{emerald!70!black}}{\thesubsubsection.}{0.5em}{}
\titlespacing*{\section}{0pt}{18pt}{8pt}

% Fancy Header/Footer setup
\pagestyle{fancy}
\fancyhf{} % clear all header and footer fields
\fancyhead[L]{\color{emerald}\small\textbf{\nouppercase{\leftmark}}}
\fancyfoot[L]{\scriptsize\color{black!50}{\lat{AI}}-সংকলিত, আলেম-যাচাইকৃত নয়}
\fancyfoot[C]{\color{emerald!70!black}\thepage}
\renewcommand{\headrulewidth}{0.4pt}
\renewcommand{\headrule}{\hbox to\headwidth{\color{emerald}\leaders\hrule height \headrulewidth\hfill}}
\renewcommand{\footrulewidth}{0pt}

% Custom boxes
% 1. Ayat/Hadith Box
\newtcolorbox{ayatbox}[1][]{
    enhanced, breakable, 
    colback=emerald!5, 
    colframe=emerald, 
    boxrule=0.5pt, 
    arc=3pt, 
    left=10pt, right=10pt, top=10pt, bottom=10pt,
    #1
}

% 2. Quote/Translation Box
\newtcolorbox{quotebox}[1][]{
    enhanced, breakable,
    frame hidden,
    colback=lightgray,
    borderline west={3pt}{0pt}{emerald},
    left=15pt, right=10pt, top=10pt, bottom=10pt,
    sharp corners,
    #1
}

% 3. AI-generated notice box
\newtcolorbox{warnbox}[1][]{
    enhanced, breakable,
    colback=red!4,
    colframe=darkred,
    boxrule=0.8pt,
    arc=3pt,
    left=10pt, right=10pt, top=10pt, bottom=10pt,
    #1
}

% Commands
\newcommand{\arab}[1]{\foreignlanguage{arabic}{#1}}
\newcommand{\kib}[1]{\textcolor{emerald}{\textbf{#1}}}

% Latin (English) fallback: Noto Serif Bengali has NO Latin glyphs
\newfontfamily\latfont[Renderer=HarfBuzz]{Noto Serif}
\newcommand{\lat}[1]{{\latfont #1}}

\setlength{\parskip}{5pt}
\setlength{\parindent}{0pt}

% Fix for missing bullet in Noto Serif Bengali
\renewcommand{\labelitemi}{$\bullet$}



\begin{document}
\begin{center}{\Large\bfseries\color{emerald} ক্ষমা করা পাপ — \lat{01}-আয়েশা-সহজ-হিসাব-ও-নবীজির-ব্যাখ্যা}\end{center}
\vspace{1em}
\section*{আয়েশা (রাঃ)-এর প্রশ্ন — \lat{“}সহজ হিসাব\lat{”} আসলে কী?}

\subsection*{সূত্র কার্ড}

\begin{table}[h]\centering\small
\begin{tabular}{ll}
বিষয় & বিবরণ \\
\hline
\textbf{ঘটনা} & আয়েশা (রাঃ) নবীজিকে (সা.) জিজ্ঞেস করেন — \lat{“}সহজ হিসাব\lat{”} (৮৪:৮) মানে কী? নবীজি (সা.) বলেন: এটি আরয (আমলের পেশকরণ); পুঙ্খানুপুঙ্খ হিসাব মানেই শাস্তি \\
\textbf{বর্ণনাকারী} & আয়েশা (রাঃ) — ইবনে আবি মুলাইকা (রহ.) থেকে \\
\textbf{গ্রন্থ ও নম্বর} & সহীহ বুখারি ৬৫৩৬; সহীহ মুসলিম ২৮৭৬ \\
\textbf{অধ্যায়} & বুখারি: কিতাবুর রিকাক — \lat{“}যার হিসাব পুঙ্খানুপুঙ্খ নেওয়া হবে সে শাস্তি পাবে\lat{”}; মুসলিম: কিতাবুল জান্নাহ — \lat{“}হিসাব প্রমাণ\lat{”} \\
\textbf{সনদের মান} & \textbf{সহীহ} (বুখারি-মুসলিম) \\
\textbf{স্থান ও সময়} & মদিনা — নবীজির (সা.) জীবদ্দশায় \\
\hline
\end{tabular}
\end{table}

\subsection*{সংক্ষিপ্ত পটভূমি}

কিয়ামতের দিনের হিসাব নিয়ে উম্মতের মনে স্বাভাবিক উদ্বেগ ছিল। কুরআনে বলা হয়েছে — যাকে আমলনামা ডান হাতে দেওয়া হবে, \lat{“}সে তো সহজ হিসাবে হিসাব পেশ করবে\lat{”} (৮৪:৭-৮)। আয়েশা (রাঃ) — যিনি কুরআন ও দ্বীন বোঝার গভীরতা জন্য বিখ্যাত — নবীজির (সা.) কাছে এই আয়াতের প্রকৃত অর্থ জিজ্ঞেস করলেন।

\subsection*{ঘটনার পূর্ণ বিবরণ}

আয়েশা (রাঃ) বলেন — রাসূলুল্লাহ (সা.) বলেছেন:

\begin{quote}
\textbf{\lat{“}কিয়ামতের দিন যার হিসাব পুঙ্খানুপুঙ্খভাবে নেওয়া হবে, সে শাস্তি পাবে।\lat{”}}
\end{quote}

আয়েশা (রাঃ) বললেন: \lat{“}আল্লাহ কি বলেননি — \textbf{"সে তো সহজ হিসাবে হিসাব পেশ করবে"} (৮৪:৮)?\lat{”}

নবীজি (সা.) বললেন:

\begin{quote}
\textbf{\lat{“}সেটি তো আরয — (আমলের পেশকরণ)। যার হিসাব পুঙ্খানুপুঙ্খভাবে নেওয়া হবে, সে শাস্তি পাবে।\lat{”}}
\end{quote}

(সহীহ বুখারি ৬৫৩৬; সহীহ মুসলিম ২৮৭৬)

\subsection*{শব্দের অর্থ}

\begin{table}[h]\centering\small
\begin{tabular}{ll}
শব্দ & অর্থ \\
\hline
\textbf{নুকিশা আল-হিসাব (\arab{نوقش} \arab{الحساب})} & পুঙ্খানুপুঙ্খ জিজ্ঞাসাবাদ — প্রতিটি কাজের খুঁটিনাটি তোলা \\
\textbf{উয্যিবা (\arab{عُذِّبَ})} & শাস্তি পাবে — ধ্বংস হবে \\
\textbf{আল-আরয (\arab{العَرْض})} & পেশকরণ — আমল দেখানো; খুঁটিনাটি জিজ্ঞাসা নয় \\
\hline
\end{tabular}
\end{table}

\subsection*{সংশ্লিষ্ট দলিল}

\begin{itemize}
\item \textbf{হাদিস (বুখারি ২৪৪১; মুসলিম ২৭৬৮):} নাজওয়া হাদিস — মুমিনকে কাছে টেনে পাপ একান্তে স্মরণ করিয়ে ক্ষমা; তারপর নেকির আমলনামা। এই হাদিসের সাথে আয়েশা (রাঃ)-এর হাদিস মিলিয়ে বোঝা যায় — মুমিনের পাপের \lat{“}মেনশন\lat{”} থাকলেও তা শাস্তি বা অপমানের জন্য নয়।
\item \textbf{আয়াত (২৫:৭০):} তাওবার পর পাপ নেকিতে বদল — হিসাবের বোঝা থেকে মুক্তি।
\item \textbf{আয়াত (১১:১১৪):} ভালো কাজ মন্দ কাজ মুছে দেয়।
\end{itemize}
\subsection*{শিক্ষা}

\begin{enumerate}
\item \textbf{\lat{“}সহজ হিসাব\lat{”} মানে হিসাবই নেই তা নয়} — বরং আমলের পেশকরণ (আরয) হবে; পুঙ্খানুপুঙ্খ জিজ্ঞাসাবাদ হবে না।
\item \textbf{ক্ষমা করা পাপ আর পুঙ্খানুপুঙ্খ জিজ্ঞাসাবাদের বিষয় নয়} — তাওবা, নেক আমল ও আল্লাহর রহমতের কারণে।
\item আয়েশা (রাঃ) কুরআনের আয়াত নিয়ে নবীজির (সা.) কাছে প্রশ্ন করলেন — \textbf{কুরআন বোঝার সঠিক পদ্ধতি} হলো রাসূলের (সা.) ব্যাখ্যা থেকে শেখা, নিজের ধারণা থেকে নয়।
\end{enumerate}

\end{document}
